\SourceLeaf{334}

\Needspace{4\baselineskip}Hieraus ist klar:

I. Haben diese drei Sinus nicht alle gleiche Zeichen, so ist die Aufgabe physisch unmöglich.

In diesem Falle wird also der Schnitt der beiden Kreise, die auf $b$ und $c$ mit den Winkeln $\mathfrak{B}, \mathfrak{C}$ beschrieben sind, nicht in den vorgeschriebenen Abschnitten liegen, und folglich in keinem von beiden, wenn
\[ \sin(\mathfrak{B}-B), \qquad \sin(\mathfrak{C}-C) \]
unter sich gleiche Zeichen haben, aber
\[ \sin(\mathfrak{A}-A) \]
das entgegengesetzte. Man wird also anstatt $\mathfrak{B}$, $\mathfrak{B}+180^0$ und anstatt $\mathfrak{C}$, $\mathfrak{C}+180^0$ an dem Schnittpunkte vorfinden.

II. Jene Bedingung ist aber zur Möglichkeit zureichend. Diess wird so bewiesen:

1) Sind die drei Sinus positiv, so fällt der Schnitt der Kreise auf $b$ und $c$ innerhalb des Kreises durch 1, 2, 3.

\textls[120]{Beweis.} Gesetzt, er fiele ausserhalb, so müsste daselbst erscheinen $\mathfrak{B}' = \mathfrak{B}+180^0$ anstatt $\mathfrak{B}$ und $\mathfrak{C}' = \mathfrak{C}+180^0$ anstatt $\mathfrak{C}$ (damit $\sin(\mathfrak{B}'-B)$ negativ werde und eben so $\sin(\mathfrak{C}'-C)$). Es bleibt aber daselbst der Winkel auf $a$
\[ =360^0-\mathfrak{B}'-\mathfrak{C}'=360^0-\mathfrak{B}-\mathfrak{C}=\mathfrak{A}. \]

2) Eben so wird für den Fall, wo die drei Sinus negativ sind, die Absurdität der Voraussetzung, der Schnitt zweier Hülfskreise falle innerhalb des Hauptkreises, bewiesen.

Die Pointe ist, dass $\mathfrak{A}, \mathfrak{B}, \mathfrak{C}$ nothwendig die Summe $=360^0$ geben müssen (d. i. durch $360^0$ divisibel). Die vorgegebenen Werthe haben diese Eigenschaft von selbst. Sind solche unverträglich, so müssen \textls[120]{zwei} abgeändert werden (um $180^0$), damit jene Eigenschaft conservirt bleibe.

Sind $b, c$ die beiden Seiten, auf denen die Hülfskreise errichtet werden, so werden die Bedingungen sich so aussprechen lassen:
\[ \sin(\mathfrak{B}-B), \quad \sin(\mathfrak{C}-C) \quad \text{und} \quad \sin(\mathfrak{B}-B+\mathfrak{C}-C) \]
sollen gleiche Zeichen haben.
